\FloatBarrier
\subsection{BPMN}\label{sec:bpmn}

Business Process Model and Notation (BPMN) is a widely-used graphical notation for
specifying business processes~\cite{BusinessProcessModel}, standardised as ISO/IEC 19510:2013.
BPMN is thus a first-class modelling language in its own right.
In the process-mining context the core elements are activities, XOR gateways (exclusive
choice), AND gateways (parallel split/join), OR gateways (inclusive choice/merge), and
sequence flows. We type BPMN's control flow directly, each sequence flow and gateway carrying an
object type, and give the resulting object-centric BPMN (OCBPMN) its cospan-algebra presentation, one further
instance of the general framework. Earlier object-centric treatments of BPMN take a different route
(Section~\ref{subsec:bpmn-related}).
Read at the level of models, this generalises the Kalenkova conversion between BPMN and Petri nets~\cite{kalenkovaProcessMiningUsing2017} to a single framework that covers four notations at once.

\begin{definition}[BPMN model]\label{def:bpmn}
  A BPMN model is a tuple
  \[
    B = (N, F, \ell, \mathcal{O}),
  \]
  where $N$ is a finite set of flow objects and $\mathcal{O}$ is a finite set of
  object types.
  Here $F \subseteq N \times \mathcal{O} \times N$ is the set of typed sequence flows,
  a typed multigraph in which an arc $(x,\omega,y)\in F$ records that an object of type
  $\omega\in\mathcal{O}$ flows from $x$ to $y$, mirroring the typed dependency multigraph
  $D$ of Definition~\ref{def:causal-net}, and
  $\ell : N \to \{\mathrm{activity}, \mathrm{event}, \mathrm{AND}, \mathrm{XOR}, \mathrm{OR}\}$
  assigns a BPMN node type to each flow object. A start event has no incoming flow and an end event
  has no outgoing flow.
  Write $A = \ell^{-1}(\mathrm{activity})$ for the activities and $G = N \setminus A$ for the
  gateways and events. Each flow carries an object type $\omega$, which is the edge typing the
  general framework uses. Gateways carry no
  object type of their own and are control-flow routing constructs
  (Remark~\ref{rem:bpmn-typed-gateways}).
  The control-flow graph of $B$ is the directed graph $(N, F)$.
\end{definition}

BPMN's OR gateways (inclusive choice) are the one construct without a direct AND/XOR
reading, and the least settled part of the standard's token semantics. The difficulty is the
OR-join, which must fire once exactly those incoming branches the matching split activated
have completed, in general a non-local decision that operational token semantics settle only
with global lookahead~\cite{dijkmanSemanticsAnalysisBusiness2008}. We sidestep the
operational question denotationally. Our concern is the structural (trace) language of a
model, what it \emph{allows} to happen, not whether a particular token run deadlocks, and
read this way an OR gateway is just the finite family of branch subsets it may activate,
each subset an ordinary AND-binding. The join then consumes exactly the subset the split
committed, so no global state is read, which is the reading object-centric causal nets and
Petri nets already use for inclusive choice.

\begin{definition}[OR-gateway semantics]\label{def:or-semantics}
  We read an OR gateway denotationally, as the set of branch subsets it allows.
  For a $k$-ary OR gateway with branch set $S$, write
  $\mathcal{P}^+(S) = \{ U \subseteq S \mid U \neq \emptyset \}$ for its $2^k - 1$
  non-empty subsets. The gateway denotes the inclusive-choice family indexed by
  $\mathcal{P}^+(S)$. An OR-split commits one $U \in \mathcal{P}^+(S)$ and activates exactly the
  branches in $U$ concurrently, and the matching OR-join fires once exactly the branches in $U$ have
  completed, the committed $U$ recorded as the gateway's binding. Each $U$ is thus an ordinary
  AND-binding and the gateway is their XOR. It is the object-centric causal-net activity with
  output-binding set $\mathrm{Out}(g)=\mathcal{P}^+(S)$, and in OCPN form the $2^k-1$ silent AND-transitions
  selected by an XOR.
\end{definition}

Because the committed $U$ is recorded as a binding, the OR-join reads exactly what the split wrote
and consults no global state. Replacing each OR gateway by its $\mathcal{P}^+(S)$ alternatives
(Definition~\ref{def:or-semantics}), each an XOR-selected AND-binding, leaves a finite AND/OR graph whose mediators are all tagged AND or XOR (Definition~\ref{def:lm-graph}), with the same signature. From there the construction is the
OCPN and OCCN one unchanged. The AND/OR graph is walked, and each activity's
gateway-mediated neighbourhoods are read off it.

This instantiates the general framework of Section~\ref{sec:general-framework}, and the two mapping
choices mirror the earlier sections. An activity is a generator. A gateway is a mediator, an AND
gateway an AND, an XOR gateway an XOR, and an OR gateway the XOR-of-ANDs above. A $1$-ary gateway is
the degree-$(1,1)$ mediator where AND, XOR and a SEQ pass-through all coincide, and simply forwards
its one branch. The gateways carry no generator of their own, exactly as the operators of the
process tree do not (Remark~\ref{rem:bpmn-typed-gateways}). Given
$B=(N,F,\ell,\mathcal{O})$, the instance map is as follows.
\medskip
\begin{center}
  \small\begin{tabularx}{\linewidth}{@{}lX@{}}
    \hline
    General framework & BPMN \\
    \hline
    AND/OR graph $(V,E)$ & the control-flow graph $(N,F)$ \\
    Activities $\mathcal{A}$ & tasks $A$ \\
    AND mediators & AND gateways, and the start and end events, where a walk ends at $\bot$ or $\top$ \\
    XOR mediators & XOR gateways, and OR gateways (exploded to XOR-of-ANDs over $\mathcal{P}^+(S)$) \\
    SEQ pass-throughs & $1$-ary gateways, where AND, XOR and SEQ coincide \\
    \hline
  \end{tabularx}
\end{center}
\medskip
The signature comes from $B$ as given. Inserting a type-preserving gateway of one input and one
output on any flow leaves every reached family, every context and hence every generator unchanged,
since a $1$-ary mediator forwards under both the AND and the XOR rule and wires are named by the
activities at their ends. The construction is therefore invariant under the normalisation of
\citet{dijkmanSemanticsAnalysisBusiness2008} without requiring it, which
matters because that normalisation is stated for their Petri-net mapping and leaves the OR-join
aside. Nothing here depends on resolving an OR-join. The signature enumerates the paths a model
permits and says nothing about how a running system would choose between them.

The engine of Section~\ref{sec:general-framework} runs unchanged, with one instance-specific
fact. Reachability is through the gateways only, an AND gateway joining one reached set from each
of its branches, since they occur together, and an XOR gateway keeping the reached sets of its
branches separate, since exactly one of them occurs. The context, generator, and composition steps
are those of Definitions~\ref{def:contexts-general} and~\ref{def:generator-cospan-general} on the
typed flows.

\begin{figure*}[tp]
  \centering
  \begin{tikzpicture}[scale=0.95, inline,
    evt/.style={circle, draw=black!70, fill=black!6, minimum size=7mm, font=\small},
    task/.style={rectangle, draw=black!65, fill=black!4, minimum size=7.5mm, font=\small},
    gw/.style={diamond, draw=black!60, fill=black!6, inner sep=1pt, minimum size=7mm, font=\small}]
  \node[evt]  (g1)  at (0,0)      {$\bot$};
  \node[gw]   (ag1) at (2.0,0)    {$+$};
  \node[gw]   (xg1) at (3.8,1.1)  {$\times$};
  \node[task] (b)   at (5.6,1.8)  {$b$};
  \node[task] (e)   at (5.6,0.5)  {$e$};
  \node[gw]   (xg2) at (7.4,1.1)  {$\times$};
  \node[task] (c)   at (5.6,-1.4) {$c$};
  \node[gw]   (ag2) at (9.2,0)    {$+$};
  \node[task] (s)   at (10.9,0)   {$s$};
  \node[evt]  (g2)  at (12.6,0)   {$\top$};

  \draw[blueflow] ([yshift=2.8pt]g1.east) -- ([yshift=2.8pt]ag1.west);
  \draw[redflow]  ([yshift=-2.8pt]g1.east) -- ([yshift=-2.8pt]ag1.west);
  \draw[blueflow] (ag1) -- (xg1);
  \draw[redflow]  (ag1) -- (c);
  \draw[blueflow] (xg1) -- (b);
  \draw[blueflow] (xg1) -- (e);
  \draw[blueflow] (b)   -- (xg2);
  \draw[blueflow] (e)   -- (xg2);
  \draw[blueflow] (xg2) -- (ag2);
  \draw[redflow]  (c)   -- (ag2);
  \draw[blueflow] ([yshift=2.8pt]ag2.east) -- ([yshift=2.8pt]s.west);
  \draw[redflow]  ([yshift=-2.8pt]ag2.east) -- ([yshift=-2.8pt]s.west);
  \draw[blueflow] ([yshift=2.8pt]s.east) -- ([yshift=2.8pt]g2.west);
  \draw[redflow]  ([yshift=-2.8pt]s.east) -- ([yshift=-2.8pt]g2.west);
\end{tikzpicture}
\unskip
  \caption{A small object-centric BPMN model in canonical form, over order (blue)
    and item (red). The start emits an order and an item, and an AND-split gateway
    ($+$) routes them down parallel branches. An XOR gateway ($\times$) verifies
    ($b$) or expedites ($e$) the order while $c$ packs the item, and an AND-join gateway
    ($+$) synchronises them for shipping ($s$). Every gateway preserves object
    types, order flows staying order-typed and item flows item-typed.}
  \label{fig:re-obpmn}
\end{figure*}

\begin{figure*}[tp]
  \centering
  \begin{tikzpicture}[inline,
    bspd/.style={spider, fill=bluecol},
    rspd/.style={spider, fill=redcol},
    mbox/.style={draw=black, rounded corners=4pt, fill=white, minimum size=0.75cm, font=\small},
    plabel/.style={font=\tiny}]

  \node[font=\footnotesize] at (2.7,1.4) {decorated cospan};
  \node[font=\footnotesize] at (9.0,1.4) {string diagram};

  \foreach \row/\g in {0/b, -1.5/e}{%
    \begin{scope}[yshift=\row cm]
      \node[font=\small] at (-2.3,0) {$g_{\g}$};
      \node[cospanblue, minimum height=0.6cm] (L) at (0,0) {};
      \node[bspd] at (0,0) {};
      \node at (0.9,0) {$\rightarrow$};
      \node[cospangrey, minimum height=1.0cm, minimum width=2.6cm] (A) at (2.7,0) {};
      \node[mbox] (f) at (2.7,0) {$\g$};
      \node[bspd] (ai) at (1.65,0) {};
      \node[bspd] (ao) at (3.75,0) {};
      \draw[bluewire] (ai) -- (f.west);
      \draw[bluewire] (f.east) -- (ao);
      \node at (4.5,0) {$\leftarrow$};
      \node[cospanblue, minimum height=0.6cm] (R) at (5.4,0) {};
      \node[bspd] at (5.4,0) {};
      \node[plabel, anchor=east] at ([shift={(-3pt,0cm)}]L.west) {$(\bot,\mathrm{o},\g)$};
      \node[plabel, anchor=west] at ([shift={(3pt,0cm)}]R.east) {$(\g,\mathrm{o},s)$};
      \node at (7.4,0) {$\mapsto$};
      \node[morphismblue, minimum size=0.75cm] (sd) at (9.0,0) {$\g$};
      \draw[bluewire] (sd.west) -- ++(-0.55,0);
      \draw[bluewire] (sd.east) -- ++(0.55,0);
    \end{scope}}

  \begin{scope}[yshift=-3.0cm]
    \node[font=\small] at (-2.3,0) {$g_c$};
    \node[cospanblue, minimum height=0.6cm] (cL) at (0,0) {};
    \node[rspd] at (0,0) {};
    \node at (0.9,0) {$\rightarrow$};
    \node[cospangrey, minimum height=1.0cm, minimum width=2.6cm] (cA) at (2.7,0) {};
    \node[mbox] (cf) at (2.7,0) {$c$};
    \node[rspd] (cai) at (1.65,0) {};
    \node[rspd] (cao) at (3.75,0) {};
    \draw[redwire] (cai) -- (cf.west);
    \draw[redwire] (cf.east) -- (cao);
    \node at (4.5,0) {$\leftarrow$};
    \node[cospanblue, minimum height=0.6cm] (cR) at (5.4,0) {};
    \node[rspd] at (5.4,0) {};
    \node[plabel, anchor=east] at ([shift={(-3pt,0cm)}]cL.west) {$(\bot,\mathrm{it},c)$};
    \node[plabel, anchor=west] at ([shift={(3pt,0cm)}]cR.east) {$(c,\mathrm{it},s)$};
    \node at (7.4,0) {$\mapsto$};
    \node[morphismblue, minimum size=0.75cm] (csd) at (9.0,0) {$c$};
    \draw[redwire] (csd.west) -- ++(-0.55,0);
    \draw[redwire] (csd.east) -- ++(0.55,0);
  \end{scope}

  \begin{scope}[yshift=-4.5cm]
    \node[font=\small] at (-2.3,0) {$g_s$};
    \node[cospanblue, minimum height=0.9cm] (sL) at (0,0) {};
    \node[bspd] at (0,0.20) {};
    \node[rspd] at (0,-0.20) {};
    \node at (0.9,0) {$\rightarrow$};
    \node[cospangrey, minimum height=1.0cm, minimum width=2.6cm] (sA) at (2.7,0) {};
    \node[mbox, minimum height=0.8cm] (sf) at (2.7,0) {$s$};
    \node[bspd] (sao) at (1.65,0.20) {};
    \node[rspd] (sai) at (1.65,-0.20) {};
    \node[bspd] (sbo) at (3.75,0.20) {};
    \node[rspd] (sbi) at (3.75,-0.20) {};
    \draw[bluewire] (sao) -- (sf.160);
    \draw[redwire]  (sai) -- (sf.200);
    \draw[bluewire] (sf.20)  -- (sbo);
    \draw[redwire]  (sf.-20) -- (sbi);
    \node at (4.5,0) {$\leftarrow$};
    \node[cospanblue, minimum height=0.9cm] (sR) at (5.4,0) {};
    \node[bspd] at (5.4,0.20) {};
    \node[rspd] at (5.4,-0.20) {};
    \node[plabel, anchor=east] at ([shift={(-3pt,0.20cm)}]sL.west) {$(b,\mathrm{o},s)$};
    \node[plabel, anchor=east] at ([shift={(-3pt,-0.20cm)}]sL.west) {$(c,\mathrm{it},s)$};
    \node[plabel, anchor=west] at ([shift={(3pt,0.20cm)}]sR.east) {$(s,\mathrm{o},\top)$};
    \node[plabel, anchor=west] at ([shift={(3pt,-0.20cm)}]sR.east) {$(s,\mathrm{it},\top)$};
    \node at (7.4,0) {$\mapsto$};
    \node[morphismblue, minimum height=0.95cm] (ssd) at (9.0,0) {$s$};
    \draw[bluewire] ($(ssd.south west)!0.70!(ssd.north west)$) -- ++(-0.55,0);
    \draw[redwire]  ($(ssd.south west)!0.30!(ssd.north west)$) -- ++(-0.55,0);
    \draw[bluewire] ($(ssd.south east)!0.70!(ssd.north east)$) -- ++(0.55,0);
    \draw[redwire]  ($(ssd.south east)!0.30!(ssd.north east)$) -- ++(0.55,0);
    \node[plabel, black!60] at (2.7,-0.72) {order from $b$ (or $e$)};
  \end{scope}
\end{tikzpicture}
\unskip
  \caption{The generator cospans of Fig.~\ref{fig:re-obpmn}. Gateways are
    mediators and yield no generators, while the activities do. The XOR gateway gives one
    generator per branch ($g_b$, $g_e$). The item thread is $g_c$, and $g_s$ is the
    object-centric synchronisation, its order port arriving from $b$ or $e$ (the
    XOR resolved upstream), so $s$ has two contexts.}
  \label{fig:re-obpmn-cospans}
\end{figure*}

\begin{exmp}[A small OCBPMN and its generator cospans]
  \label{ex:running-bpmn-cospan}
  Consider the OCBPMN of Fig.~\ref{fig:re-obpmn}, over the object types
  $\mathrm{order}$ and $\mathrm{item}$. Every split and join is an explicit gateway,
  and each gateway preserves object types. The AND-split routes the order and item
  by type, the XOR offers two same-type alternatives for the order, and nothing
  creates or converts a type (Remark~\ref{rem:bpmn-typed-gateways}).

  The gateways are mediators, so only the activities yield generators
  (Fig.~\ref{fig:re-obpmn-cospans}). The XOR gateway gives one generator per branch,
  $g_b:\{(\bot,\mathrm{o},b)\}\to\{(b,\mathrm{o},s)\}$ and
  $g_e:\{(\bot,\mathrm{o},e)\}\to\{(e,\mathrm{o},s)\}$. The item thread is
  $g_c:\{(\bot,\mathrm{it},c)\}\to\{(c,\mathrm{it},s)\}$. The AND-join is
  the object-centric synchronisation, shown in its $b$-context,
  \begin{gather*}
    g_s = \Bigl(\{(b,\mathrm{o},s),(c,\mathrm{it},s)\}
      \xrightarrow{\iota} A \xleftarrow{\iota}\twocolbreak 
      \{(s,\mathrm{o},\top),(s,\mathrm{it},\top)\},\; d_s\Bigr).
  \end{gather*}
  Its left boundary carries both object types. The order arrives from $b$ or from $e$ (the XOR
  resolved upstream), so $s$ has two contexts, the $e$-context replacing $b$ by $e$. The AND
  gateways correspond to $s$'s AND-binding and the XOR gateway to the order's alternatives, matching
  the causal-net bindings of Section~\ref{sec:causal-nets}. Any count constraint (a key-distribution
  split, say) relates the leg counts of $s$, is recorded in its constraint system $\Lambda$
  (Definition~\ref{def:multiplicity-decoration}) and needs no further machinery.
\end{exmp}

\begin{figure*}[tp]
  \centering
  \resizebox{\textwidth}{!}{\begin{tikzpicture}[inline,
    evt/.style={circle, draw=black!70, fill=black!6, minimum size=6.5mm, font=\small},
    task/.style={rectangle, draw=black!65, fill=black!4, minimum size=7mm, font=\small},
    gw/.style={diamond, draw=black!60, fill=black!6, inner sep=0.5pt, minimum size=7mm,
               font=\scriptsize},
    bspd/.style={spider, fill=bluecol},
    mbox/.style={draw=black, rounded corners=4pt, fill=white, minimum size=0.75cm, font=\small},
    plabel/.style={font=\tiny}]

  \node[font=\footnotesize, anchor=west] at (-1.6,1.9) {(a) OCBPMN OR gadget};
  \node[evt]  (g1) at (0,0)     {$\bot$};
  \node[gw]   (os) at (1.5,0)   {OR};
  \node[task] (v)  at (3.0,0.7) {$v$};
  \node[task] (u)  at (3.0,-0.7){$u$};
  \node[gw]   (oj) at (4.5,0)   {OR};
  \node[task] (s)  at (6.0,0)   {$s$};
  \draw[blueflow] (g1) -- (os);
  \draw[blueflow] (os) -- (v);  \draw[blueflow] (os) -- (u);
  \draw[blueflow] (v)  -- (oj);  \draw[blueflow] (u) -- (oj);
  \draw[blueflow] (oj) -- (s);
  \node[font=\footnotesize, anchor=west] at (7.1,0.35)
    {denotes the XOR over $\mathcal{P}^+(\{v,u\})$};
  \node[font=\footnotesize, anchor=west] at (7.1,-0.35)
    {$=\{\{v\},\{u\},\{v,u\}\}$ (last is $\mathrm{AND}(v,u)$)};

  \node[font=\footnotesize, anchor=west] at (-1.6,-1.5) {(b) the three contexts of $s$};
  \begin{scope}[xshift=1.0cm]
  \node[font=\footnotesize] at (2.7,-2.1) {decorated cospan};
  \node[font=\footnotesize] at (9.0,-2.1) {string diagram};
  \foreach \row/\lab/\src in {-3.1/{U=\{v\}}/v, -4.6/{U=\{u\}}/u}{%
    \begin{scope}[yshift=\row cm]
      \node[font=\scriptsize] at (-2.3,0) {$\lab$};
      \node[cospanblue, minimum height=0.6cm] (L) at (0,0) {};
      \node[bspd] at (0,0) {};
      \node at (0.9,0) {$\rightarrow$};
      \node[cospangrey, minimum height=1.0cm, minimum width=2.6cm] (A) at (2.7,0) {};
      \node[mbox] (f) at (2.7,0) {$s$};
      \node[bspd] (ai) at (1.65,0) {};
      \node[bspd] (ao) at (3.75,0) {};
      \draw[bluewire] (ai) -- (f.west);
      \draw[bluewire] (f.east) -- (ao);
      \node at (4.5,0) {$\leftarrow$};
      \node[cospanblue, minimum height=0.6cm] (R) at (5.4,0) {};
      \node[bspd] at (5.4,0) {};
      \node[plabel, anchor=east] at ([shift={(-3pt,0cm)}]L.west) {$(\src,\mathrm{o},s)$};
      \node[plabel, anchor=west] at ([shift={(3pt,0cm)}]R.east) {$(s,\mathrm{o},\top)$};
      \node at (7.4,0) {$\mapsto$};
      \node[morphismblue, minimum size=0.75cm] (sd) at (9.0,0) {$s$};
      \draw[bluewire] (sd.west) -- ++(-0.55,0);
      \draw[bluewire] (sd.east) -- ++(0.55,0);
    \end{scope}}
  \begin{scope}[yshift=-6.1cm]
    \node[font=\scriptsize] at (-2.3,0) {$U=\{v,u\}$};
    \node[cospanblue, minimum height=0.9cm] (L) at (0,0) {};
    \node[bspd] at (0,0.20) {};
    \node[bspd] at (0,-0.20) {};
    \node at (0.9,0) {$\rightarrow$};
    \node[cospangrey, minimum height=1.0cm, minimum width=2.6cm] (A) at (2.7,0) {};
    \node[mbox, minimum height=0.8cm] (f) at (2.7,0) {$s$};
    \node[bspd] (ai1) at (1.65,0.20) {};
    \node[bspd] (ai2) at (1.65,-0.20) {};
    \node[bspd] (ao) at (3.75,0) {};
    \draw[bluewire] (ai1) -- (f.160);
    \draw[bluewire] (ai2) -- (f.200);
    \draw[bluewire] (f.east) -- (ao);
    \node at (4.5,0) {$\leftarrow$};
    \node[cospanblue, minimum height=0.6cm] (R) at (5.4,0) {};
    \node[bspd] at (5.4,0) {};
    \node[plabel, anchor=east] at ([shift={(-3pt,0.20cm)}]L.west) {$(v,\mathrm{o},s)$};
    \node[plabel, anchor=east] at ([shift={(-3pt,-0.20cm)}]L.west) {$(u,\mathrm{o},s)$};
    \node[plabel, anchor=west] at ([shift={(3pt,0cm)}]R.east) {$(s,\mathrm{o},\top)$};
    \node at (7.4,0) {$\mapsto$};
    \node[morphismblue, minimum height=0.95cm] (sd) at (9.0,0) {$s$};
    \draw[bluewire] ($(sd.south west)!0.70!(sd.north west)$) -- ++(-0.55,0);
    \draw[bluewire] ($(sd.south west)!0.30!(sd.north west)$) -- ++(-0.55,0);
    \draw[bluewire] (sd.east) -- ++(0.55,0);
  \end{scope}
  \end{scope}
\end{tikzpicture}
\unskip}
  \caption{(a) An OR gadget over $S=\{v,u\}$, order-typed. It denotes the XOR over the non-empty
    subsets $\mathcal{P}^+(S)=\{\{v\},\{u\},\{v,u\}\}$, the last an AND-binding. (b) The resulting
    three contexts of $s$, each a decorated cospan and the string diagram it denotes. The
    $\{v,u\}$ context gives $s$ two order ports.}
  \label{fig:re-bpmn-or}
\end{figure*}

\begin{exmp}[An OR gateway and its alternatives]
  \label{ex:bpmn-or}
  Suppose an order must undergo at least one of two optional checks, verify
  ($v$) and audit ($u$), before ship ($s$). An OR-split over the branch
  set $S=\{v,u\}$ feeds an OR-join into $s$ (Fig.~\ref{fig:re-bpmn-or}). By
  Definition~\ref{def:or-semantics} the gateway denotes the family indexed by
  $\mathcal{P}^+(S)=\bigl\{\,\{v\},\{u\},\{v,u\}\,\bigr\}$, so $s$ gains one context per committed
  subset, with left boundaries
  \[
    \{(v,\mathrm{o},s)\},\quad \{(u,\mathrm{o},s)\},\quad
    \{(v,\mathrm{o},s),(u,\mathrm{o},s)\},
  \]
  the last being the AND-binding in which both checks were taken. The branches
  contribute $g_v:\{(\bot,\mathrm{o},v)\}\to\{(v,\mathrm{o},s)\}$ and, likewise,
  $g_u$. The inclusive choice is thus an XOR over these three AND-bindings, exactly the
  $\mathcal{P}^+(S)$ explosion, and the OR-join consumes precisely the binding the split
  committed, so no global token state is consulted.
\end{exmp}

\begin{remark}[Typed gateways are routing-only]
  \label{rem:bpmn-typed-gateways}
  Every gateway is type-preserving, carrying the same set of object types on its incoming and
  outgoing flows and serving only to route objects. A gateway is a mediator, so it has no port of
  its own and the walk of Definition~\ref{def:contexts-general} passes through it. A gateway may not
  convert one type into another. Any count constraint on the routing (a key-distribution split, a
  batch cardinality) relates the leg counts of an activity it routes to and enters that activity's
  constraint system $\Lambda$ (Definition~\ref{def:multiplicity-decoration}) without further
  machinery.
\end{remark}

\begin{theorem}[Canonical cospan-algebra presentation of BPMN]\label{thm:bpmn-cospan}
  Let $B$ be any finite BPMN model.
  Then $B$ canonically determines a cospan-algebra signature $\Sigma$, and hence the free
  symmetric monoidal category $\mathbf{F}(\Sigma)$.
\end{theorem}

\begin{proof}
  The delta is the inclusive choice. By Definition~\ref{def:or-semantics} each OR gateway is the
  $\mathcal{P}^+(S)$ family of XOR-selected AND-bindings, a finite replacement that changes no
  signature, and afterwards every mediator is an AND or an XOR, a $1$-ary gateway being the
  degree-$(1,1)$ case, and an edge joining two activities directly is the degenerate case,
  traversed in one step (Definition~\ref{def:lm-graph}), and inserting a type-preserving $1$-ary
  gateway on such an edge changes no reached family and so no generator. The schema of Section~\ref{sec:notation-by-notation} then applies at the instance
  map above, the tasks being the activities, each edge carrying its flow type $\omega$
  (Definition~\ref{def:bpmn}), and finiteness of $N$ giving the conditions of
  Definition~\ref{def:lm-graph}. The mediators are the same AND/XOR kinds as the Petri- and
  causal-net instances, and the reach families are finite even under cyclic gateways, since the
  graph is finite and a branch stops at any node already on its path
  (Definition~\ref{def:contexts-general}).
\end{proof}

The presentation needs no structural precondition on $B$, and the OR-join, which operational token semantics settle only with global lookahead~\cite{dijkmanSemanticsAnalysisBusiness2008}, is handled denotationally.

